\documentclass[11pt]{article}
\usepackage{mathblatt}
\makeatletter
\setlength{\mbpfbe}{0mm}% keine Punkte (Bauregel 6.4)
\makeatother
% Rechenraum als Karo (Bauregel 6.9): n Rechenschritte = n mal 10 mm
\newcommand{\karo}[1]{\par\smallskip\noindent\begin{tikzpicture}
  \clip (0,0) rectangle ({\linewidth-1mm},{#1*10mm});
  \draw[black!16,line width=0.3pt,step=5mm] (0,0) grid ({\linewidth},{#1*10mm});
  \end{tikzpicture}\par}
\newcommand{\linie}[1]{\,{\color{black!45}\rule[-1pt]{#1}{0.4pt}}\,}
\newcommand{\pkt}{\textperiodcentered{}}
\begin{document}

\pfheftstil
\pflaufkopf[P10 \pkt{} mehr]{Flächeninhalt und Umfang}{PG40-7}\pfseitevon
\pfheftkopf{Flächeninhalt und Umfang}{mehr}
\pfstufenzeile{Formeln umstellen \pkt{} Formeln \pkt{} Fläche gesucht}
% 2020-OS-B1d-m1 (Schwere I, rang 12)
\begin{pfaufg}{nach~P10\\’20~\pkt{}~1d}{1.}{}
Ein Quadrat hat den Flächeninhalt 81 cm². Gib seine Seitenlänge an.
\karo{2}
\fusshilfe{\mbox{1:~9 cm}}
\end{pfaufg}
% 2019-OS-B1e-m1 (Schwere I, rang 13)
\begin{pfaufg}{nach~P10\\’19~\pkt{}~1e}{2.}{}
Gib eine Gleichung für den Flächeninhalt A der gesamten Fläche an.
\par\smallskip\noindent{\small\itshape Skizze: Viereck (Rechteck): liegend, durch eine senkrechte Linie geteilt; linker Teil grau mit Breite a, rechter Teil kariert mit Breite b; gemeinsame Höhe c rechts}\par
\hfill A $=$\linie{30mm}
\fusshilfe{\mbox{2:~(a + b) · c}}
\end{pfaufg}
% 2018-OS-B1f-m1 (Schwere I, rang 14)
\begin{pfaufg}{nach~P10\\’18~\pkt{}~1f}{3.}{}
Ein Rechteck hat den Umfang u = 30 cm; die Seite a ist 11 cm lang. Gib die Länge der Seite b an.
\par\smallskip\noindent{\small\itshape Skizze: Viereck (Rechteck): liegend, unten a = 11 cm, rechts b; nicht maßstabsgerecht}\par
\karo{2}
\fusshilfe{\mbox{3:~4 cm}}
\end{pfaufg}
% 2017-OS-K3a-m1 (Schwere I, rang 15)
\begin{pfaufg}{nach~P10\\’17~\pkt{}~3a}{4.}{}
Familie Sommer hat einen Swimmingpool gebaut. Seine Grundfläche ist ein Rechteck, an dessen kurze Seiten je ein Halbkreis angesetzt ist. Die kurze Seite des Rechtecks ist a = 6 m lang, die lange Seite b = 11 m. Gib an, aus welchen Teilflächen die Grundfläche des Pools besteht.
\par\smallskip\noindent{\small\itshape Skizze: Figur aus dem Vorspann // zusammengesetzte Figur: Stadionform, Rechteck mit zwei angesetzten Halbkreisen, schraffiert; Breite a = 6 m links, gerades Stück b = 11 m unten; nicht maßstabsgerecht}\par
\hfill Ein Rechteck und zwei Halbkreise (zusammen ein Kreis mit Durchmesser a) $=$\linie{30mm}
\fusshilfe{\mbox{4:~Ein Rechteck und zwei Halbkreise (zusammen ein Kreis mit Durchmesser a)}}
\end{pfaufg}
% 2016-OS-K3b-m1 (Schwere I, rang 16)
\begin{pfaufg}{nach~P10\\’16~\pkt{}~3b}{5.}{}
Beim Kugelstoßen steht man in einem Abwurfring; das ist ein Kreis mit 2,30 m Durchmesser. Ermittle den Flächeninhalt dieses Kreises.
\karo{2}
\fusshilfe{\mbox{5:~4,15 m²}}
\end{pfaufg}
% 2014-OS-B1g-m1 (Schwere I, rang 17)
\begin{pfaufg}{nach~P10\\’14~\pkt{}~1g}{6.}{}
Die graue Figur ist aus gleich großen Quadraten mit der Seitenlänge a zusammengesetzt. Gib einen Term für ihren Umfang u an.
\par\smallskip\noindent{\small\itshape Skizze: Kästchenfigur: graues Kreuz auf Kästchenraster (Kästchen = a); waagerechter Balken 7 Kästchen breit und 1 hoch, senkrechter Balken 1 breit und 5 hoch, beide mittig gekreuzt; linke Kante mit a beschriftet}\par
\karo{2}
\fusshilfe{\mbox{6:~24 · a}}
\end{pfaufg}
% 2020-OS-K7b-m1 (Schwere II, rang 18)
\begin{pfaufg}{nach~P10\\’20~\pkt{}~7b}{7.}{}
Im Dreieck ABC ist bei C ein rechter Winkel; BC ist 15 m und AC 28 m lang. Der Punkt Q liegt auf AC. Das Dreieck QBC hat den Flächeninhalt 120 m². Berechne die Länge der Strecke QC.
\par\smallskip\noindent{\small\itshape Skizze: rechtwinkliges Dreieck: rechter Winkel bei C (oben links), B rechts von C, A senkrecht unter C; Q auf AC nahe C, QB gestrichelt; nicht maßstabsgerecht}\par
\karo{3}
\fusshilfe{\mbox{7:~16 m}}
\end{pfaufg}
% 2019-OS-K4c-m1 (Schwere II, rang 19)
\begin{pfaufg}{nach~P10\\’19~\pkt{}~4c}{8.}{}
Auf einem Mast sitzt ein Prisma, dessen Grund- und Deckfläche gleichseitige Dreiecke mit 1,80 m Seitenlänge sind; die Höhe dieser Dreiecke beträgt 1,56 m. Das Prisma ist 1,50 m hoch. Ermittle den Flächeninhalt der Deckfläche. Durch ein Loch ist Regenwasser in das Prisma gelaufen. Zeige mit einer Rechnung, dass etwa 2,1 m³ Wasser hineinpassen. Leer wiegt das Prisma 250 kg; 1 m³ Regenwasser wiegt 1 000 kg. Berechne, wie schwer das Prisma ist, wenn ein Zehntel seines Innenraums (10 \%) voll Wasser ist.
\par\smallskip\noindent{\small\itshape Skizze: Dreieck: gleichseitig, grau, mit gestrichelter Höhe $h_{\text{Deckfläche}}$ und rechtem Winkel; nicht maßstabsgerecht}\par
\karo{4}
\fusshilfe{\mbox{8:~1,40 m²; 2,1 m³; 461 kg}}
\end{pfaufg}
% 2018-OS-K6a-m1 (Schwere II, rang 20)
\begin{pfaufg}{nach~P10\\’18~\pkt{}~6a}{9.}{}
Ein runder Turm mit 5,8 m Durchmesser steht auf einer Rasenfläche, die 18,0 m × 11,0 m groß ist. Der Rasen um den Turm herum soll neu angelegt werden. Berechne die Fläche des neuen Rasens.
\par\smallskip\noindent{\small\itshape Skizze: zusammengesetzte Figur: Rechteck (gepunktet) mit grauem Kreis darin; nicht maßstabsgerecht}\par
\karo{2}
\fusshilfe{\mbox{9:~171,6 m²}}
\end{pfaufg}
% 2017-OS-K3b-m1 (Schwere II, rang 21)
\begin{pfaufg}{nach~P10\\’17~\pkt{}~3b}{10.}{}
Familie Sommer hat einen Swimmingpool gebaut. Seine Grundfläche ist ein Rechteck, an dessen kurze Seiten je ein Halbkreis angesetzt ist. Die kurze Seite des Rechtecks ist a = 6 m lang, die lange Seite b = 11 m. Berechne den Flächeninhalt der Grundfläche des Pools.
\par\smallskip\noindent{\small\itshape Skizze: Figur aus dem Vorspann // zusammengesetzte Figur: Stadionform, Rechteck mit zwei angesetzten Halbkreisen, schraffiert; Breite a = 6 m links, gerades Stück b = 11 m unten; nicht maßstabsgerecht}\par
\karo{2}
\fusshilfe{\mbox{10:~94,3 m²}}
\end{pfaufg}
% 2015-OS-K6c-m1 (Schwere II, rang 22)
\begin{pfaufg}{nach~P10\\’15~\pkt{}~6c}{11.}{}
Mia baut das Modell einer Pyramide mit quadratischer Grundfläche: Die Grundkante ist 0,90 m lang, die Pyramide ist 0,70 m hoch. Die Seitenflächen sägt sie aus dünnem Sperrholz. Weise nach, dass sie für eine Seitenfläche etwa 0,37 m² Sperrholz braucht.
\karo{3}
\fusshilfe{\mbox{11:~0,83 m; 0,37 m²}}
\end{pfaufg}
% 2014-OS-K5c-m1 (Schwere II, rang 23)
\begin{pfaufg}{nach~P10\\’14~\pkt{}~5c}{12.}{}
Eine Kühlvitrine hat die Form eines Prismas. Ihre Grundfläche ist ein gleichschenkliges Trapez mit den parallelen Seiten 108 cm und 84 cm und dem Flächeninhalt 3 360 cm². Die Höhe h dieses Trapezes ist die Tiefe der Vitrine. Weise nach, dass Kuchenplatten mit 30 cm Durchmesser in die Vitrine passen.
\par\smallskip\noindent{\small\itshape Skizze: Viereck (Trapez): parallele Seiten 108 cm (oben) und 84 cm (unten), Schenkel 37 cm, Höhe h mit rechten Winkeln; nicht maßstabsgerecht}\par
\karo{2}
\fusshilfe{\mbox{12:~35 cm}}
\end{pfaufg}
% 2014-OS-K5d-m1 (Schwere II, rang 24)
\begin{pfaufg}{nach~P10\\’14~\pkt{}~5d}{13.}{}
Die Rückwand derselben Vitrine ist ein Rechteck, 108 cm breit und 150 cm hoch. Sie soll aus einer rechteckigen Platte gesägt werden, die 1,25 m breit ist und einen Flächeninhalt von 1,8 m² hat. Entscheide mit einer Rechnung, ob die Platte dafür reicht.
\karo{2}
\fusshilfe{\mbox{13:~1,44 m hoch}}
\end{pfaufg}
% 2015-OS-K5d-m1 (Schwere III, rang 25)
\begin{pfaufg}{nach~P10\\’15~\pkt{}~5d}{14.}{}
Im Dreieck AED ist bei E ein rechter Winkel. Die Hypotenuse AD ist 3,9 cm lang, der Winkel bei D beträgt $\delta$ = 28°. Der Flächeninhalt des Dreiecks AED soll berechnet werden. Gib an, welche Seiten und Winkel du dazu brauchst. Schreibe die Rechenschritte in einer sinnvollen Reihenfolge auf.
\par\smallskip\noindent{\small\itshape Skizze: rechtwinkliges Dreieck: rechter Winkel bei E, Hypotenuse AD = 3,9 cm, $\delta$ = 28° bei D; Teil des Drachenvierecks ABCD (Vorspann 2015-OS-K5 in wortlaut-eigen-dreiecke.csv); nicht maßstabsgerecht}\par
\karo{5}
\fusshilfe{\mbox{14:~28°; 1,83 cm; 3,44 cm; 3,15 cm²}}
\end{pfaufg}
\par\vfill\noindent{\scriptsize\color{mbgrau}Originale: P10 2020 \pkt{} 1d, P10 2019 \pkt{} 1e, P10 2018 \pkt{} 1f, P10 2017 \pkt{} 3a, P10 2016 \pkt{} 3b, P10 2014 \pkt{} 1g, P10 2020 \pkt{} 7b, P10 2019 \pkt{} 4c, P10 2018 \pkt{} 6a, P10 2017 \pkt{} 3b, P10 2015 \pkt{} 6c, P10 2014 \pkt{} 5c, P10 2014 \pkt{} 5d, P10 2015 \pkt{} 5d\par}
\end{document}